% ECONOMICS_PROBLEM: {"id": "C73-113", "title": "Constant regret in general-sum Markov games", "jel_code": "C73", "source_id": "OP-0151"}
% !TeX program = xelatex
\documentclass[11pt,a4paper]{article}
\usepackage[margin=23mm]{geometry}
\usepackage{fontspec}
\setmainfont{Latin Modern Roman}
\usepackage{amsmath,amssymb,mathtools}
\usepackage{xcolor,enumitem,booktabs,longtable,array,xurl,hyperref}
\definecolor{muted}{HTML}{53636C}
\hypersetup{colorlinks=true,urlcolor=blue,linkcolor=blue}
\setlength{\parindent}{0pt}
\setlength{\parskip}{5pt}
\newcommand{\R}{\mathbb R}
\newcommand{\E}{\mathbb E}
\newcommand{\N}{\mathbb N}
\newcommand{\ind}{\mathbf 1}
\newcommand{\Var}{\operatorname{Var}}
\newcommand{\Cov}{\operatorname{Cov}}
\newcommand{\supp}{\operatorname{supp}}
\DeclareMathOperator*{\argmin}{arg\,min}
\DeclareMathOperator*{\argmax}{arg\,max}
\newcommand{\entrymeta}[3]{{\sffamily\small\color{muted}\textbf{\texttt{#1}}\quad|\quad #2\par #3\par}}
\newcommand{\Problemref}[1]{\texttt{#1}}
\newcommand{\JELref}[2]{\texttt{#2} (JEL \texttt{#1})}
\begin{document}
\section*{Constant regret in general-sum Markov games}
\textbf{Problem ID:} \texttt{C73-113}\quad \textbf{JEL:} \texttt{C73}\par
% BEGIN ECONOMICS STATEMENT
\label{problem:OP-0151}
% GAME_THEORY_PROBLEM: {"local_id": "GT-021", "original_number": 21, "kind": "R", "entry_kind": "statement_only", "published": false, "review_required": true, "category": "game_theory"}
\label{game:GT-021}
{\sffamily\small\color{muted}
{\sffamily\small\color{muted}\textbf{GT-021} | Learning in games\par R | Precise formulation of a source research agenda; source-date open\par}
\par}
\subsection*{Definitions and mathematical statement}
A finite episodic Markov game has players $N=\{1,\ldots,n\}$, states $S$, nonempty finite action sets $A_i$, horizon $H\geq1$, initial state $s_1$, transition kernels $P_h(\cdot\mid s,a)$, and rewards $r_{i,h}(s,a)\in[0,1]$. All realized states and actions are observed. A behavioral policy assigns an action distribution to each player-observed history within an episode; $\Pi_i$ denotes this set. For a profile $\pi$, put
\[
 V_i(\pi)=\mathbb E_{\pi}\!\left[\sum_{h=1}^{H}r_{i,h}(s_h,a_h)\right].
\]
For profiles $\pi^1,\ldots,\pi^T$, define episode-wise external regret against a single policy by
\[
 R_i(T)=\sup_{\tau_i\in\Pi_i}\sum_{t=1}^{T}
 \bigl(V_i(\tau_i,\pi^t_{-i})-V_i(\pi^t)\bigr).
\]
An editorial version of the constant-regret target is a decentralized self-play procedure and a finite constant $C_G$, independent of $T$, such that
\[
 \forall G\quad\forall T\geq1\quad\forall i\in N,\qquad R_i(T)\leq C_G.
\]
Use the source's full-feedback regime: at a queried $(h,s,a_i)$ the learner receives the exact own reward and transition kernel averaged against the opponents' current Markov policies. A computational formulation must additionally fix a per-round oracle budget and an arithmetic-operation budget, both polynomial in the finite game description and $T$; these resource restrictions are editorial specifications.

Let $\mu_T$ select one joint profile $\pi^t$ uniformly, before the episode. A coarse deviation ignores its own profile recommendation and uses an arbitrary $\tau_i\in\Pi_i$. Linearity gives
\[
 \max_i\sup_{\tau_i\in\Pi_i}
 \bigl(V_i(\tau_i,(\mu_T)_{-i})-V_i(\mu_T)\bigr)
 =\frac1T\max_i R_i(T).
\]
Thus the stated regret target implies a $C_G/T$ coarse-correlated-equilibrium gap for this particular mixture.
\subsection*{Short English statement}
Can decentralized self-play achieve cumulative external regret bounded independently of the number of episodes, with a corresponding $1/T$ equilibrium error?
\subsection*{Variants and scope boundaries}
The quantifier is over opponents generated by the same self-play procedure. Constant regret against arbitrary adversarial opponents is a different target. The source's recursive correlated-policy output and its local weighted regrets must not be identified with the ordinary empirical mixture above without an additional argument. Sample-based learning with unknown rewards or transitions changes the feedback model.
\subsection*{Source relationship and status}
[S1], Section 5, explicitly retains constant regret as a future question after an $O(\log T/T)$ CCE guarantee. The displayed external-regret problem is an explicit editorial specialization; the source does not completely specify the resource model for every possible constant-regret algorithm.
\subsection*{References}
\begingroup\footnotesize\raggedright
\textbf{[S1]} Asr\i n Efe Yorulmaz and Tamer Ba\c{s}ar, \emph{Near Optimal Convergence to Coarse Correlated Equilibrium in General-Sum Markov Games}, arXiv:2511.02157v1 (2025). Locators: Sections 2.1--2.5 (game, oracle, policy, CCE and regret definitions), Section 5 (constant-regret question). \url{https://arxiv.org/html/2511.02157v1}
\par\endgroup

\subsection*{Common conventions: Source mathematical conventions}

\textbf{Finite games.} For a finite set $A$, $\Delta(A)$ is its probability
simplex. Players choose independent mixed strategies in Nash equilibrium;
correlation is allowed only when explicitly part of the solution concept.
In a game with payoff functions $u_i$ and strategy profile $x$, an additive
$\varepsilon$ Nash equilibrium satisfies
\[
 u_i(x)\geq u_i(x_i',x_{-i})-\varepsilon
 \quad\text{for every player $i$ and every admissible deviation $x_i'$}.
\]
For numerical approximation, payoffs are normalized as locally specified.
An $\varepsilon$ payoff residual is distinct from distance at most
$\varepsilon$ to the set of exact equilibria. Point convergence, convergence
of empirical play, and convergence in distance to an equilibrium set are
also distinct.

\textbf{Computational representations.} Unless stated otherwise, a complexity
question takes rational data with binary encoding; polynomial time is measured
in total bit length. A fixed-accuracy polynomial algorithm is not necessarily
polynomial in $1/\varepsilon$, and neither is necessarily polynomial in
$\log(1/\varepsilon)$. Succinct games and value, demand, or sampling oracles
must declare their interfaces. Exact real solutions require an explicit
representation; an unspecified list of arbitrary real numbers is not a
finite computational output. A claimed algorithmic barrier is meaningful
only for its specified model and reductions.

\textbf{Dynamic games.} Histories, information, observation, and allowable
randomization are part of the model. A stationary strategy depends only on
the current state; a positional strategy in a deterministic graph game
chooses one action at each controlled vertex. Nash equilibrium at an initial
state and equilibrium after every history are different requirements.
An expectation of a pathwise $\liminf$ payoff, a $\liminf$ of expected averages,
a limiting expected average, and a uniform finite-horizon equilibrium are
different criteria. None is silently substituted for another.

\textbf{Allocation and mechanisms.} A complete allocation assigns every item
exactly once unless otherwise stated. Zero-valued goods, negative values,
capacity constraints, feasibility, lotteries, and copies of goods affect
fairness definitions and are specified locally. Dominant-strategy incentive
compatibility, Bayesian incentive compatibility, universal truthfulness,
and truthfulness in expectation impose different inequalities. Whenever
an objective ratio has zero denominator, its convention is stated or those
instances are excluded.

\textbf{Infinite-dimensional models.} Expectations, integrals, stochastic
equations, and optimization objectives require their local regularity,
integrability, filtrations, and admissible domains. A backward--forward
mean-field system is a boundary-value problem; it is not an initial-value
semiflow without an additional construction. Long-time, large-population,
and vanishing-noise limits require an order of limits and a convergence
topology. If a minimum need not be attained, the objective is its infimum
or supremum and the approximation requirement is stated separately.
% END ECONOMICS STATEMENT
\end{document}
